\begin{frame}[label=q3c1]{3c.1 · Decimales: menores por departamento}
  \small
  \pregunta{¿Qué porcentaje del total de cada departamento corresponde a menores de 18 años, y cuál es el control nacional?}
  \entrada{RDD; $d$ = departamento, $c$ = cantidad. Menores: 0--5, 6--11 y 12--17 años.}
  \begin{columns}[T,onlytextwidth]
    \column{0.55\textwidth}
    \textbf{Map y Reduce de cada operación}
    \parmr{A: total}{$(d,c)$\\Todas las filas}
      {\texttt{reduceByKey(\_+\_)}\\$(d,T_d)$: total por dpto.}
    \parmr{B: menores}{$(d,c)$\\Filtrar menores antes del Map}
      {\texttt{reduceByKey(\_+\_)}\\$(d,M_d)$: menores por dpto.}
    \parmr[Proyección]{C: control}{\texttt{values} de A y B\\Tomar los totales parciales}
      {\texttt{reduce(\_+\_)}\\Total y menores nacionales}
    \column{0.42\textwidth}
    \textbf{Salida}\\[5pt]
    \footnotesize
    \begin{tabular}{@{}lrr@{}}
      \toprule Departamento & Menores/total & \%\\\midrule
      Madre de Dios & 473/655 & 72.21\\
      Ucayali & 1\,504/2\,126 & 70.74\\
      Lima Met. & 19\,751/39\,490 & 50.02\\\bottomrule
    \end{tabular}
    \par\vspace{6pt}\textbf{Cálculo después de reducir}\\
    Unir A y B por departamento:\\
    $p_d=100M_d/T_d$.
    \par\vspace{5pt}Control: $70\,178/121\,427=57.79\%$.
    \par\vspace{2pt}
    \codigo{rdd/03c_1_menores_por_departamento}{20}{20}
  \end{columns}
  \respuesta{Madre de Dios: 72.21\% de menores; el denominador es el total reportado del departamento.}
  \origen{rdd}{03c_1_menores_por_departamento}
  \note{40 s. Leer cada fila de izquierda a derecha: primero la clave y el valor del Map, luego su Reduce. C proyecta values y reduce los parciales. Después del join se calcula la división.}
\end{frame}
